\label{seminar05}
\subsection{Цель и подготовка}

В этой лабораторной две части, выполняемые \textbf{по порядку}. Сначала вы
создадите локальный репозиторий Git и зафиксируете исходную версию двух
программ. В том же репозитории проверите игнорирование файлов, ветки и
обмен с другим локальным клоном. Затем соберёте программы в Linux и
исследуете коды возврата, процессы и их память. После опытов измените
один исходник и зафиксируете результат ещё одним коммитом.

\subsection{Порядок сдачи лабораторной работы}

\begin{enumerate}
    \item На занятии выполните часть работы, зафиксируйте команды и наблюдения и покажите результат преподавателю. Если в процессе работы возникли вопросы, задайте их преподавателю.
    \item К следующему лабораторному занятию завершите работу, изучите использованные команды и понятия, подготовьтесь объяснить работу программ, механизмов Git, процессов и памяти. Подготовьте отчёт по обеим частям согласно разделу «Отчёт» и отправьте его преподавателю в формате PDF или HTML. Сохраните исходники, оба рабочих каталога и репозиторий отдельно для демонстрации.
    \item На следующем занятии индивидуально защитите работу: покажите преподавателю отчёт и результат, объясните свои действия и ответьте на вопросы. Преподаватель сообщит, если нужно что-то исправить.
    \item На последнем занятии защитите предыдущую лабораторную работу и сдайте оставшиеся задолженности. Преподаватель объявит текущую накопленную оценку за лабораторные работы.
\end{enumerate}

Для работы нужны терминал Linux, \texttt{git}, \texttt{gcc}, \texttt{ps},
\texttt{file} и два файла \texttt{status\_demo.c},
\texttt{isolation\_demo.c} из комплекта материалов семинара №5. Команды ниже
рассчитаны на Bash в Debian/Ubuntu; знак приглашения \texttt{\$} не вводите.
Проверьте инструменты:

\begin{verbatim}
git --version
gcc --version
ps --version
\end{verbatim}

Если программ нет, в учебной Debian/Ubuntu установите пакеты (могут
потребоваться права администратора):

\begin{verbatim}
sudo apt update
sudo apt install git gcc libc6-dev file procps
\end{verbatim}

Создайте отдельный рабочий каталог:

\begin{verbatim}
mkdir -p ~/aiya/lab05
\end{verbatim}

Поместите в него три выданных файла: \texttt{README.md},
\texttt{status\_demo.c} и \texttt{isolation\_demo.c}. Затем перейдите
в рабочий каталог и проверьте его содержимое:

\begin{verbatim}
cd ~/aiya/lab05
ls
\end{verbatim}

Перед переходом к части I убедитесь, что \texttt{ls} показывает все три
файла. Создавайте Git-репозиторий только в своём рабочем каталоге.

\subsection{Часть I. Первый локальный репозиторий Git}

Git хранит историю \emph{файлов проекта}.
\textbf{Рабочая директория} содержит текущие файлы, \textbf{индекс} (staging
area) --- то, что войдёт в следующий коммит, \textbf{коммит} фиксирует
снимок проекта. Удалённый сервис и учётная запись для этой работы не нужны.

\subsubsection*{Создайте и проверьте репозиторий}

\begin{verbatim}
git init
git status
\end{verbatim}

Ожидайте сообщения о новом репозитории и списке ещё не отслеживаемых файлов
(untracked). Название исходной ветки может зависеть от настроек Git. Если
Git при коммите попросит имя автора, настройте \emph{только этот} репозиторий:
\texttt{git config -{}-local user.name "Ваше имя"} и
\texttt{git config -{}-local user.email "ваш адрес"}. Не копируйте в отчёт
чужие имя или адрес из примера.

Исполняемые файлы и временные результаты не следует добавлять в историю.
Откройте \texttt{nano .gitignore} и запишите строки (сохранить: Ctrl+O,
Enter; выйти: Ctrl+X):

\begin{verbatim}
/status_demo
/isolation_demo
/input.txt
/result.txt
/error.txt
*.o
\end{verbatim}

Начальный символ \texttt{/} ограничивает правило корнем этого репозитория.
Проверьте, что хотите сохранить именно исходники и инструкции, затем
зафиксируйте первую версию:

\begin{verbatim}
git add README.md status_demo.c isolation_demo.c .gitignore
git status
git diff --cached
git commit -m "Add starting files for lab 5"
git log -1 --oneline
git status
\end{verbatim}

\texttt{git diff -{}-cached} показывает изменения в индексе до коммита;
обычный \texttt{git diff} показывает ещё не добавленные изменения рабочего
файла. Если Git открыл длинный вывод в просмотрщике \texttt{less}, нажмите
\textbf{q} для выхода. После коммита \texttt{git status} должен сообщать об отсутствии
незафиксированных изменений. Запишите в отчёт краткий идентификатор первого
коммита и своими словами объясните, что попало в него. Команда
\mbox{\texttt{git add .}} добавляет все подходящие изменения из текущего каталога;
в этой работе используйте имена файлов, чтобы перед коммитом видеть
точный состав индекса. \texttt{git config -{}-global} настраивает автора
для всех репозиториев пользователя, а применённое здесь
\texttt{-{}-local} --- только для учебного каталога.

\subsubsection*{Игнорирование и восстановление}

Создайте два пробных файла и попробуйте добавить один из них в индекс:

\begin{verbatim}
printf 'temporary\n' > scratch.txt
printf 'debug\n' > debug.log
git add scratch.txt
git status --short
git rm --cached scratch.txt
git status --short
\end{verbatim}

\texttt{git rm -{}-cached} удалил \texttt{scratch.txt} из \emph{индекса},
но оставил файл на диске: он снова не отслеживается. Добавьте в уже
созданный \texttt{.gitignore} две строки \texttt{/scratch.txt} и
\texttt{*.log}. Убедитесь, что Git их игнорирует:

\begin{verbatim}
git check-ignore -v scratch.txt debug.log
git status --short
\end{verbatim}

Файл \texttt{.gitignore} сам по себе отслеживается, поэтому его изменение
появится в статусе. Правило игнорирования не удаляет файл с диска и не
перестаёт отслеживать уже \emph{закоммиченный} файл автоматически.
Теперь внесите пробную правку в отслеживаемый README и отмените её:

\begin{verbatim}
printf 'Trial line\n' >> README.md
git diff -- README.md
git restore README.md
git status --short
\end{verbatim}

\texttt{git restore} без дополнительных флагов возвращает рабочую версию
README к состоянию индекса: пробная строка пропадёт. Старый синтаксис
\texttt{git checkout -{}- README.md} решает сходную задачу, но
\texttt{git switch} удобнее для перехода между ветками. Не восстанавливайте
этой командой изменения, которые хотите сохранить.

\subsubsection*{Ветка и слияние}

Ветка (branch) позволяет фиксировать отдельную линию изменений. Имя
начальной ветки зависит от настроек Git, поэтому сохраните его в
переменной Bash, затем создайте ветку для краткой заметки:

\begin{verbatim}
base_branch=$(git branch --show-current)
git switch -c feature-notes
git branch
printf 'Git and processes\n' > notes.txt
git add notes.txt .gitignore
git diff --cached
git commit -m "Add notes and ignore temporary files"
git switch "$base_branch"
git merge --no-ff feature-notes -m "Merge feature notes"
git branch -d feature-notes
git log --graph --oneline --all -6
\end{verbatim}

До слияния в исходной ветке не было \texttt{notes.txt}; после него файл
появился. \texttt{git switch -c} создаёт ветку и переходит на неё,
\texttt{git merge} переносит изменения, \texttt{git branch -d} удаляет
уже влитую ветку. Флаг \texttt{-{}-no-ff} в примере сохраняет отдельный
коммит слияния. Альтернатива \texttt{git switch} --- \texttt{git checkout}.

\subsubsection*{Локальный удалённый репозиторий}

Git умеет обмениваться коммитами без сайта и сети. Создайте рядом с
\texttt{lab05} \emph{bare-репозиторий} --- место хранения истории без
рабочих исходников. Для этого вернитесь в первый терминал и убедитесь,
что находитесь в \texttt{\textasciitilde{}/aiya/lab05}:

\begin{verbatim}
pwd
remote_dir="$HOME/aiya/lab05-remote.git"
copy_dir="$HOME/aiya/lab05-copy"
git init --bare "$remote_dir"
git remote add origin "$remote_dir"
git push -u origin "$base_branch"
git clone -b "$base_branch" "$remote_dir" "$copy_dir"
git -C "$copy_dir" log -1 --oneline
\end{verbatim}

\texttt{origin} --- имя связи с соседним репозиторием; это учебная модель
удалённого репозитория, хотя оба каталога находятся на одной машине.
\texttt{clone} создал независимую рабочую копию. Флаг \texttt{-b}
выбирает сохранённую ветку даже тогда, когда настройка начальной ветки
bare-репозитория отличается. При повторном выполнении используйте новые
имена каталогов: Git не клонирует поверх существующего.

В клоне сделайте и отправьте небольшую правку. Автор для этого второго
рабочего каталога настраивается отдельно, если он не задан глобально:

\begin{verbatim}
git -C "$copy_dir" config --local user.name "Ваше имя"
git -C "$copy_dir" config --local user.email "ваш адрес"
printf 'Copied work\n' >> "$copy_dir/notes.txt"
git -C "$copy_dir" add notes.txt
git -C "$copy_dir" commit -m "Update note in clone"
git -C "$copy_dir" push origin "$base_branch"
git pull --ff-only origin "$base_branch"
cat notes.txt
\end{verbatim}

Последние две команды выполняются в \emph{первом} каталоге
\texttt{lab05}. Вторая строка появилась в нём после \texttt{pull}; флаг
\texttt{-{}-ff-only} разрешает только прямое продвижение ветки. В отчёте
объясните, где теперь лежат два рабочих каталога и история \texttt{origin}.

\subsubsection*{Дополнительно: спрятать или отменить изменение}

Следующие опыты не нужны для основного маршрута. Выполняйте их после
обязательных шагов и \emph{не переписывайте историю основного каталога}.
Для \texttt{stash} можно временно убрать незакоммиченную правку README
и вернуть её:

\begin{verbatim}
printf 'Temporary note\n' >> README.md
git stash push -m "Temporary note" -- README.md
git status --short
git stash pop
git restore README.md
\end{verbatim}

Чтобы различить два способа отмены коммита, сначала сделайте отдельную
копию для опытов (из \texttt{lab05}):

\begin{verbatim}
recovery_dir="$HOME/aiya/lab05-recovery"
git clone . "$recovery_dir"
git -C "$recovery_dir" config --local user.name "Ваше имя"
git -C "$recovery_dir" config --local user.email "ваш адрес"
printf 'Trial\n' >> "$recovery_dir/notes.txt"
git -C "$recovery_dir" add notes.txt
git -C "$recovery_dir" commit -m "Trial note"
git -C "$recovery_dir" reset HEAD~1
git -C "$recovery_dir" status --short
\end{verbatim}

\texttt{reset HEAD\textasciitilde{}1} передвигает ветку назад, но
сохраняет пробную правку в рабочем файле. Сделайте из неё новый коммит
и отмените его \texttt{revert}, который \emph{добавит} ещё один коммит:

\begin{verbatim}
git -C "$recovery_dir" add notes.txt
git -C "$recovery_dir" commit -m "Trial again"
git -C "$recovery_dir" revert --no-edit HEAD
git -C "$recovery_dir" log -3 --oneline
\end{verbatim}

Команда \texttt{git reset -{}-hard} также стирает несохранённые изменения;
в опыте она не требуется. \texttt{reset} меняет историю текущей ветки,
а \texttt{revert} сохраняет прежние коммиты и создаёт новый.

\subsection{Часть II. Коды возврата, процессы и изоляция}

\subsubsection*{Соберите и запустите программы}

\begin{verbatim}
gcc -std=c11 -Wall -Wextra -O0 status_demo.c -o status_demo
gcc -std=c11 -Wall -Wextra -O0 isolation_demo.c -o isolation_demo
file status_demo isolation_demo
git status --short
\end{verbatim}

\texttt{.c} --- исходный файл, \texttt{status\_demo} и
\texttt{isolation\_demo} --- исполняемые файлы на диске. При запуске
возникает процесс; после его завершения файл остаётся. \texttt{file}
поможет увидеть тип файлов. Сборка создаёт ELF-файлы для вашей архитектуры;
они не должны появиться в \texttt{git status -{}-short} благодаря
\texttt{.gitignore}.

\subsubsection*{Исследуйте код возврата и потоки}

Каждую пару команд выполняйте подряд: \texttt{\$?} в Bash означает статус
\emph{последней} завершённой команды и меняется после следующей команды.

\begin{verbatim}
./status_demo --code=0
echo $?
./status_demo --code=7
echo $?
./status_demo --code=oops
echo $?
\end{verbatim}

В первом случае ожидается 0, во втором 7, при неверном аргументе --- 2 и
сообщение в \texttt{stderr}. Для этой учебной программы допустимые явные
коды 0--125: число 0 принято считать успешным завершением, ненулевое
указывает на иной результат; смысл конкретного ненулевого кода задаёт
программа. Исследуйте сохранение обычного вывода и диагностики отдельно:

\begin{verbatim}
./status_demo --code=7 > result.txt
echo $?
cat result.txt
./status_demo --code=oops 2> error.txt
echo $?
cat error.txt
\end{verbatim}

В первом случае перенаправление вывода не меняет код 7; сообщение об ошибке
во втором случае окажется в \texttt{error.txt}. После \texttt{cat} значение
\texttt{\$?} уже относится к \texttt{cat}. Дополнительно проверьте ввод из
файла и канал на готовых командах:

\begin{verbatim}
printf 'alpha\nbeta\n' > input.txt
cat < input.txt | wc -l
echo $?
\end{verbatim}

Ожидается две строки. Статус конвейера Bash без специальных настроек ---
обычно статус \emph{последней} команды, здесь \texttt{wc}. Файлы
\texttt{input.txt}, \texttt{result.txt},
\texttt{error.txt} остаются только для опытов.

Чтобы увидеть право исполнения, сначала запишите ожидаемый результат, затем
проверьте и верните разрешение:

\begin{verbatim}
chmod -x status_demo
./status_demo --code=0
echo $?
chmod +x status_demo
./status_demo --code=0
echo $?
\end{verbatim}

После снятия бита исполнения оболочка обычно сообщает об отказе в доступе;
её статус ошибки не является кодом, который вернула \texttt{main} программы:
она не была запущена. После восстановления права программа снова вернёт 0.
Аналогично проверьте попытки запустить несуществующий путь
\texttt{./no\_such\_program} и имя \texttt{status\_demo} без
\texttt{./} (если текущий каталог не входит в \texttt{PATH}). Запишите
сообщения оболочки и объясните отличие \texttt{no such file or directory}
и \texttt{command not found} от кода, возвращённого самой C-программой.

\subsubsection*{Родитель, потомок и отдельные значения}

\begin{verbatim}
./isolation_demo
echo $?
\end{verbatim}

Программа вызывает \texttt{fork()}, создавая дочерний процесс. Оба
печатают свои PID, числовое значение адреса локальной переменной и её
значение. После запуска ребёнок меняет \texttt{value} с 41 на 99 и
посылает родителю сигнал через pipe; родитель выводит значение \emph{после}
изменения ребёнком. Ожидайте: разные PID, одинаковое число виртуального
адреса в этом конкретном запуске, 99 у ребёнка и по-прежнему 41 у родителя.
PID и само число адреса каждый раз могут меняться. \texttt{waitpid()}
позволяет родителю узнать код ребёнка 7. При этом сама запущенная из Bash
программа возвращает \texttt{0}: \texttt{echo \$?} показывает \emph{её},
а не непосредственно дочерний код.

Откройте \textbf{второй терминал}. В первом запустите
\texttt{./isolation\_demo -{}-observe}: ребёнок останется работающим около
8 секунд. Сразу перепишите напечатанный \texttt{child pid} и во втором
терминале выполните (подставьте своё число):

\begin{verbatim}
ps -p 12345 -o pid,ppid,stat,comm
\end{verbatim}

Ожидается строка ребёнка с его PID и PID родителя в столбце PPID; если
строки уже нет, повторите опыт и выполните \texttt{ps} быстрее. Не
используйте PID из этой методички буквально. Два процесса могут разделять
разрешённые ресурсы, например унаследованный открытый канал; наличие
\texttt{pipe} не означает общей переменной \texttt{value}. Опыт показывает
разделение обычной памяти родителя и потомка, \emph{но не доказывает}, что
абсолютно любой доступ одного процесса к другому невозможен: ОС управляет
правами и может разрешать общую память и специальные средства отладки.

\subsubsection*{Измените исходник и зафиксируйте результат}

Найдите в \texttt{isolation\_demo.c} присваивание ребёнку
\texttt{value = 99;} и замените \texttt{99} на \texttt{73}. Снова соберите
и запустите \texttt{isolation\_demo}; объясните, почему у родителя
по-прежнему 41. Теперь проверьте различие между рабочим файлом и индексом:

\begin{verbatim}
git diff
git add isolation_demo.c
git diff --cached
git commit -m "Change child value in isolation experiment"
git log -2 --oneline
git status --short
\end{verbatim}

В истории основного репозитория должны быть начальный коммит, коммит
ветки, коммит слияния, правка из клона и итоговое изменение C-программы;
собранные файлы и результаты запусков в неё не входят. Сохраните оба
рабочих каталога и bare-репозиторий для показа
преподавателю, не отправляйте в публичный репозиторий персональные данные.

\subsection{Отчёт}

Составьте \textbf{один отчёт} по обеим частям: цель, команды,
краткий граф истории с начальным, веточным и итоговым C-коммитами;
объяснение, что сделали \texttt{git rm -{}-cached},
\texttt{git restore}, \texttt{merge}, \texttt{clone}, \texttt{push},
\texttt{pull}; ожидаемые и фактические коды трёх запусков,
различие stdout/stderr, PID и адреса родителя и ребёнка, значения до и
после изменения, объяснение результата \texttt{ps}. Ответьте на вопросы:
почему процесс не равен исполняемому файлу? Почему одинаковое число
виртуального адреса не означает одну переменную? Чей код показывает
\texttt{echo \$?} после \texttt{./isolation\_demo}? Почему наличие
общего файлового ресурса не отменяет изоляцию обычной памяти?
